% Source: figures/edit-breakdown.dat (figures/edit-breakdown.py가 reports/experiments/main/edit_tasks.json에서 생성)
% 원본 321개에 대한 전체 편집 시간를 변경, 추가, 삭제·페이지 설정로 나눈 비율(%).
% 막대는 원자료의 소수 두 자리를 유지하고, 추가 비중은 여기에서 한 번만 정수 반올림한다.
\begin{figure}[H]
\centering
\begin{tikzpicture}
\begin{axis}[
  xbar stacked, bar width=3.1mm, scale only axis, width=0.76\linewidth, height=24mm,
  xmin=0, xmax=100, xtick={0,25,50,75,100}, xlabel={시간 비중(\%)},
  y dir=reverse, ytick=data, yticklabels from table={figures/edit-breakdown.dat}{name},
  tick label style={font=\scriptsize}, label style={font=\footnotesize}, xlabel shift=-2pt,
  enlarge y limits=0.08, axis line style={black!60}, tick style={black!60},
  legend style={font=\scriptsize, at={(0.5,1.02)}, anchor=south, legend columns=3, draw=none,
    /tikz/every even column/.append style={column sep=6pt}},
  legend image code/.code={\draw[#1] (0cm,-0.1cm) rectangle (0.3cm,0.1cm);}]
\addplot[fill=black!22, draw=black!60] table[x=edit, y expr=\coordindex] {figures/edit-breakdown.dat};
% 추가 구간의 중앙에 해당 비중만 표시한다.
\addplot[fill=black!65, draw=black!65,
  nodes near coords={\pgfmathprintnumber[fixed,precision=0]{\pgfplotspointmeta}},
  every node near coord/.append style={font=\scriptsize,text=white}]
  table[x=add, y expr=\coordindex] {figures/edit-breakdown.dat};
\addplot[fill=white, draw=black!60] table[x=del, y expr=\coordindex] {figures/edit-breakdown.dat};
\legend{변경, 추가, 기타}
\end{axis}
\end{tikzpicture}
\caption{편집 유형별 시간 비중}
\label{fig:edit-breakdown}
\end{figure}
